% ============================================================
%  Section 5 — Agents
% ============================================================
\section{Agent Design}
\label{sec:agents}

Each agent loads its system prompt from a versioned Markdown file at
\texttt{skills/<company>/<agent>/vNNN.md}. The loaded version is logged in
SQLite traces, making the prompt used for a decision inspectable. Agents are
independent: updating the Extractor skill does not touch the Coder or Validator.

\subsection{LLM Client}

The LLM-calling stages use a shared \texttt{LLMClient} that wraps Azure AI
Foundry and OpenAI-compatible endpoints; the deterministic Router does not call
an LLM. Provider and model roles are configured through environment variables.
The current client configuration:
\begin{itemize}[leftmargin=*, itemsep=1pt]
  \item Sets temperature to 0 for model calls.
  \item Reads reasoning effort per role from configuration; the reported
        baseline and latency benchmark used \texttt{reasoning\_effort=none}.
  \item Removes returned reasoning content before passing text downstream;
        the client logs an approximate word count when that content is present.
  \item With \texttt{DEMO\_CACHE=1}, responses are cached by a SHA-256 hash
        of the full prompt. Cache hits incur zero network calls — useful for
        reproducible demos.
\end{itemize}

\subsection{Router}

The \agent{Router} determines the document type (\textit{invoice}, \textit{receipt},
\textit{purchase\_order}) and input modality (\textit{text} or \textit{photo}).
In our current implementation it is a deterministic Code node in n8n using file
extension and keyword matching — no LLM call. This avoids a model request but
can misroute files with unfamiliar extensions or ambiguous text; it is not
fail-proof.

\subsection{Extractor}

The \agent{Extractor} converts a document into a structured JSON object with all
fields the Validator and Coder need: supplier, date, document number, each line item
with its amounts, HTVA, TVA, TTC, Timbre Fiscal, and the matricule fiscal.
For text-layer PDFs, text is extracted via \texttt{pypdf}. For photos, the image
is sent to the vision deployment and the model's transcription is treated as the
text input.

The skill prompt uses four techniques:
\begin{enumerate}[leftmargin=*, itemsep=1pt]
  \item \textbf{Role prompting} — ``You are an expert Tunisian accounting clerk.''
  \item \textbf{Structured output schema} — the expected JSON schema is listed
        in the system prompt with field names and types.
  \item \textbf{Format guardrails} — explicit instructions for the TND amount
        format, the matricule fiscal pattern, and the date format.
  \item \textbf{Termination instruction} — ``Respond only with the JSON object,
        no explanatory text.''
\end{enumerate}

\subsection{Coder}

The \agent{Coder} assigns a ledger account code and cites the rule it used. Before
the LLM call, the Coder retrieves the five most relevant rule chunks from Qdrant
(filtered to the current company) and injects them into the system prompt. The
score is binary: 1.0 if the code matches the gold label \emph{and} the cited rule
appears as a substring in the retrieved chunks or skill file; 0.0 otherwise. This
is a mechanical support check. It does not verify that the cited text semantically
justifies the code, so the metric is weaker than a human audit of the rationale.

\begin{figure}[ht]
\centering
\begin{tikzpicture}[
  every node/.style={font=\scriptsize},
  box/.style={draw, rounded corners=3pt, minimum width=2.1cm,
              minimum height=0.6cm, thick, align=center},
  arr/.style={-Stealth, thick},
]
  \node[box, fill=C5!15, draw=C5]  (q)  at (0,0)    {Supplier + doc type\\(query)};
  \node[box, fill=C5!25, draw=C5]  (qd) at (3.2,0)  {Qdrant search\\top-5 chunks};
  \node[box, fill=C1!15, draw=C1]  (sk) at (0,-1.4) {Skill file\\v001.md};
  \node[box, fill=C1!25, draw=C1]  (lm) at (3.2,-1.4){LLM\\(1 call)};
  \node[box, fill=C3!20, draw=C3]  (out)at (3.2,-2.8){account\_code\\+ cited\_rule};

  \draw[arr] (q)  -- node[above, font=\scriptsize]{retrieve} (qd);
  \draw[arr] (qd) -- node[right, font=\tiny]{inject chunks} (lm);
  \draw[arr] (sk) -- node[below, font=\tiny]{system prompt} (lm);
  \draw[arr] (lm) -- (out);
\end{tikzpicture}
\caption{Coder agent flow: RAG retrieves relevant rules, which are injected
alongside the skill prompt into a single LLM call.}
\label{fig:coder_flow}
\end{figure}

\subsection{Validator}

The \agent{Validator} runs in two phases. First, three deterministic Python checks
run without any LLM call:
\begin{itemize}[leftmargin=*, itemsep=1pt]
  \item Arithmetic: $|\text{HTVA} + \text{TVA} - \text{TTC}| \leq 0.005$ currency units.
  \item Matricule fiscal present and non-empty.
  \item A duplicate marker check if \texttt{duplicate\_of} is present in the extracted payload.
\end{itemize}
If any check fails, a single LLM call generates a one-sentence explanation for
the human reviewer, and \texttt{needs\_review} is set to \texttt{true}.

Alongside the supplied \texttt{duplicate\_of} marker, the API registers a
normalized company, supplier, invoice-number, and date key in SQLite; a match
identifies the earlier document and routes the new item for review. Replays of
the same document ID are idempotent. Keeping validation logic in Python rather than delegating it to the LLM is a
deliberate choice: arithmetic at a 0.005 currency-unit tolerance is computed in Python and
unreliable in an LLM, especially for Tunisian formatting with three decimal places.

\subsection{Exception Handler}

When the Validator raises a flag, the \agent{Exception Handler} prepares a
structured escalation dossier: a summary of the errors, the recommended action
(correct the arithmetic / re-request the document / contact the supplier), and
the responsible contact. This is stored as a JSON object in SQLite for later
review; the current workflow does not send an external notification to a reviewer.
